\documentclass[12pt]{article}

\title{Animate\LaTeX{} for M\textsc{atlab}}			
\author{Per Bergstr\"{o}m}		
\date{March 05, 2014}					

\usepackage{animate}
\usepackage{graphicx}           

\begin{document}
\maketitle

\section{Tangent Line}

$f'(x_{0})=\displaystyle \lim_{h \rightarrow 0} \frac{f(x_{0}+h)-f(x_{0})}{h}$

\begin{center}
\input{exTangLine/exTangLine}
\end{center}

\newpage

\section{Taylor Polynomial}

(Maclaurin polynomial)

\noindent
\vspace{2.5em}$P_n(x)=f(0)+f'(0)x+\displaystyle\frac{f''(0)}{2}x^2+\ldots+\displaystyle\frac{f^{(n)}(0)}{n!}x^n$


\begin{center}
\input{exTaylorPoly/exTaylorPoly}
\end{center}

\begin{center}
\input{exTaylorPoly2/exTaylorPoly2}
\end{center}

\newpage

\section{The Newton-Raphson Method}
$f(x)=0$

\noindent
\vspace{2.5em}$x_{k+1}=x_{k}-\displaystyle\frac{f(x_{k})}{f'(x_{k})},\;k=0,1,\ldots$

\begin{center}
\input{exNewton/exNewton}
\end{center}

\begin{center}
\input{exNewton2/exNewton2}
\end{center}

\newpage
\section{Least Squares}
$\hat{\mathbf{x}} = \displaystyle \arg \min_{\mathbf{x}} || \mathbf{A}\mathbf{x}-\mathbf{b} ||_2^2$

\begin{center}
\input{exLS/exLS}
\end{center}						

\newpage

\section{Curvature}

\begin{center}
\input{exCurvature/exCurvature}
\end{center}

\section{Solid of Revolution}

\begin{center}
\input{exSolidRev/exSolidRev}

(Press the play button again after the pauses)
\end{center}

\newpage

\section{Riemann Sums}

$U(f,P)=\displaystyle \sum_{i=1}^n f(l_i)\Delta x_i$\newline
$L(f,P)=\displaystyle \sum_{i=1}^n f(u_i)\Delta x_i$

\begin{center}
\input{exRiemann/exRiemann}
\end{center}

\newpage

\section{Sine and Cosine}
Moving sine and cosine function
\begin{center}
\input{example1/example1}
\end{center}
The animation is created using \texttt{animategraphicsLaTeX}.

\section{Rotating Points}
Points on a sphere are rotating around its center
\begin{center}
\input{example2/example2}
\end{center}
The animation is created using \texttt{animateinlineLaTeX}.

\newpage

\section{Animations}
Lots of different animations can easily be created using \texttt{animateLaTeX} in \textsc{Matlab}. This might for example be appropriate in pdf-presentations. The command \texttt{$\backslash$animategraphics} is used in \texttt{animategraphicsLaTeX} and the environment \texttt{'animateinline'} is used in \texttt{animateinlineLaTeX}. Both \texttt{$\backslash$animategraphics} and \texttt{'animateinline'} are in the \LaTeX{} \texttt{animate} Package\footnote{http://ctan.uib.no/macros/latex/contrib/animate/animate.pdf}. In \texttt{animateinlineLaTeX} each frame might have different settings, such as trimming, size etc. \texttt{$\backslash$includegraphics} is used for inserting graphics into the frames. The settings are inserted into the optional command inside the brackets. See documentation for \texttt{$\backslash$includegraphics} for more infromation.

\end{document}

